\subsection{Matrices}

\begin{exercise}[3c1]
    Suppose $T \in \lin{V, W}$. Show that with respect to each choice of bases of $V$ and $W$, the matrix of $T$ has at least $\dim \range T$ nonzero entries.
\end{exercise}

\begin{sol}
    Let $v_1, \ldots, v_n$ be a basis of $V$ and $w_1, \ldots, w_m$ be a basis of $W$. Then the matrix of $T$ with respect to these bases is given by
    \[\mm(T) = 
    \begin{pmatrix}
        Tv_1 & Tv_2 & \cdots & Tv_n
    \end{pmatrix}
    \]
    where each $Tv_j$ is expressed as a linear combination of the basis vectors $w_1, \ldots, w_m$. Let $\dim \range T = r$. Then there exist $r$ linearly independent vectors in $\range T$. After possibly reordering the basis vectors of $V$, we can assume that $Tv_1, \ldots, Tv_r$ are linearly independent. Since each $Tv_j$ is a linear combination of the basis vectors of $W$, each $Tv_j$ must have at least one nonzero entry in its corresponding column of the matrix $\mm(T)$. Therefore, there are at least $r$ nonzero entries in the matrix $\mm(T)$, which shows that the matrix of $T$ has at least $\dim \range T$ nonzero entries.
\end{sol}

\begin{exercise}[3c2]
    Suppose $V$ and $W$ are finite-dimensional and $T \in \lin{V, W}$. Prove that $\dim \range T = 1$ if and only if there exist a basis of $V$ and a basis of $W$ such that with respect to these bases, all entries of $\mm(T)$ equal $1$.
\end{exercise}

\begin{sol}
    Suppose $\dim \range T = 1$. Then there exists a nonzero vector $y_1 \in W$ such that $\range T = \spn y_1$. Let $\dim V = n$ and $\dim W = m$. Extend $y_1$ to a basis $y_1, \ldots, y_m$ of $W$, and define the following vectors in $W$:
    \[w_1 = y_1 - (y_2 + \cdots + y_m), \quad w_k = y_k \text{ for } k = 2, \ldots, m.\]
    It is clear that $w_1 + w_2 + \cdots + w_m = y_1$, and hence $w_1, \ldots, w_m$ spans $W$. Now suppose for scalars $\alpha_1, \ldots, \alpha_m \in \f$, we have
    \[\alpha_1 w_1 + \alpha_2 w_2 + \cdots + \alpha_m w_m = 0.\]
    Then we can express this as
    \[\alpha_1(y_1 - (y_2 + \cdots + y_m)) + \alpha_2 y_2 + \cdots + \alpha_m y_m = 0.\]
    Simplifying, we get
    \[(\alpha_1) y_1 + (-\alpha_1 + \alpha_2) y_2 + \cdots + (-\alpha_1 + \alpha_m) y_m = 0.\]
    Since $y_1, \ldots, y_m$ are linearly independent, we must have $\alpha_1 = \alpha_2 = \cdots = \alpha_m = 0$. Therefore, $w_1, \ldots, w_m$ are linearly independent and hence form a basis of $W$.

    We now construct a basis of $V$. By the rank nullity theorem, we have
    \[\dim V = \dim \null T + \dim \range T = \dim \null T + 1.\]
    Let $u_2, \ldots, u_n$ be a basis of $\null T$. Then we can extend this to a basis $u_1, u_2, \ldots, u_n$ of $V$, where $u_1 \notin \null T$. Since $\range T = \spn y_1$, we have
    \[Tu_1 = \beta y_1\]
    for some nonzero scalar $\beta \in \f$. If $\beta \neq 1$, we may replace $u_1$ with $\beta\inv u_1$ to ensure that $Tu_1 = y_1$. We construct the following vectors in $V$:
    \[v_1 = u_1, \quad v_k = u_1 + u_k \text{ for } k = 2, \ldots, n.\]
    Note that for every $j = 1, \ldots, n$, we have $Tv_j = Tu_1 + Tu_j = y_1 + 0 = y_1$. We claim that $v_1, \ldots, v_n$ form a basis of $V$. Suppose for scalars $\alpha_1, \ldots, \alpha_n \in \f$, we have
    \[\alpha_1 v_1 + \alpha_2 v_2 + \cdots + \alpha_n v_n = 0.\]
    Then we can express this as
    \[(\alpha_1 + \alpha_2 + \cdots + \alpha_n) u_1 + \alpha_2 u_2 + \cdots + \alpha_n u_n = 0.\]
    Since $u_1, \ldots, u_n$ are linearly independent, we must have $\alpha_1 + \alpha_2 + \cdots + \alpha_n = 0$ and $\alpha_2 = \cdots = \alpha_n = 0$. This implies that $\alpha_1 = 0$, and hence $v_1, \ldots, v_n$ are linearly independent and form a basis of $V$. Then with respect to the bases $v_1, \ldots, v_n$ of $V$ and $w_1, \ldots, w_m$ of $W$, we have that $Tv_j = y_1 = w_1 + w_2 + \cdots + w_m$ for each $j = 1, \ldots, n$. Therefore, the matrix of $T$ with respect to these bases has all entries equal to $1$.

    Conversely, suppose there exist a basis of $V$ and a basis of $W$ such that with respect to these bases, all entries of $\mm(T)$ equal $1$. Then for each basis vector $v_j$ of $V$, we have
    \[Tv_j = w\]
    for some fixed nonzero vector $w \in W$, where $w = w_1 + \cdots + w_m$ is the sum of the basis vectors of $W$. Since $\range T$ is spanned by the images of the basis vectors of $V$, then any sum of $Tv_j$ is a scalar multiple of $w$. Therefore, $\range T = \spn w$, and hence $\dim \range T = 1$.
\end{sol}

\begin{exercise}[3c3]
    Suppose $v_1, \ldots, v_n$ is a basis of $V$ and $w_1, \ldots, w_m$ is a basis of $W$.
    \begin{subexercises}
        \item Show that if $S, T \in \lin{V, W}$, then $\mm(S + T) = \mm(S) + \mm(T)$.
        \item Show that if $\lambda \in \f$ and $T \in \lin{V, W}$, then $\mm(\lambda T) = \lambda \mm(T)$.
    \end{subexercises}
\end{exercise}

\begin{sole}
    \item Let $v_1, \ldots, v_n$ be a basis of $V$ and $w_1, \ldots, w_m$ be a basis of $W$. Then the matrix of $S + T$ with respect to these bases is given by
    \begin{align*}
        \mm(S + T) & = 
        \begin{pmatrix}
            (S + T)v_1 & (S + T)v_2 & \cdots & (S + T)v_n
        \end{pmatrix} \\
        & = 
        \begin{pmatrix}
            Sv_1 + Tv_1 & Sv_2 + Tv_2 & \cdots & Sv_n + Tv_n
        \end{pmatrix} \\
        & = 
        \begin{pmatrix}
            Sv_1 & Sv_2 & \cdots & Sv_n
        \end{pmatrix} +
        \begin{pmatrix}
            Tv_1 & Tv_2 & \cdots & Tv_n
        \end{pmatrix} \\
        & = \mm(S) + \mm(T).
    \end{align*}
    \item Let $v_1, \ldots, v_n$ be a basis of $V$ and $w_1, \ldots, w_m$ be a basis of $W$. Then the matrix of $\lambda T$ with respect to these bases is given by
    \begin{align*}
        \mm(\lambda T) & = 
        \begin{pmatrix}
            (\lambda T)v_1 & (\lambda T)v_2 & \cdots & (\lambda T)v_n
        \end{pmatrix} \\
        & = 
        \begin{pmatrix}
            \lambda Tv_1 & \lambda Tv_2 & \cdots & \lambda Tv_n
        \end{pmatrix} \\
        & = \lambda
        \begin{pmatrix}
            Tv_1 & Tv_2 & \cdots & Tv_n
        \end{pmatrix} \\
        & = \lambda \mm(T). \qh
    \end{align*}
\end{sole}

\begin{exercise}[3c4]
    Suppose that $D \in \lin{\pp_3(\r), \pp_2(\r)}$ is the differentiation map defined by $Dp = p'$. Find a basis of $\pp_3(\r)$ and a basis of $\pp_2(\r)$ such that the matrix of $D$ with respect to these bases is
    \[\begin{pmatrix}
        1 & 0 & 0 & 0 \\
        0 & 1 & 0 & 0 \\
        0 & 0 & 1 & 0
    \end{pmatrix}.\]
\end{exercise}

\begin{sol}
    Consider the standard basis of $\pp_3(\r)$. Reorder it to obtain the basis $x^3, x^2, x, 1$. Consider the standard basis of $\pp_2(\r)$. Reorder it to obtain the basis $x^2, x, 1$. Observe that $D(x^3) = 3x^2$, $D(x^2) = 2x$, $D(x) = 1$, and $D(1) = 0$. However, we need the matrix of $D$ to have the desired form. To achieve this, we can scale the basis vectors of $\pp_3(\r)$ and $\pp_2(\r)$ appropriately. Then we may define the basis for $\pp_3(\r)$ as
    \[\frac 13 x^3, \frac 12 x^2, x, 1\]
    while maintaining the standard basis for $\pp_2(\r)$ as $x^2, x, 1$. Then we have
    \[D\left(\frac 13 x^3\right) = x^2, \quad D\left(\frac 12 x^2\right) = x, \quad D(x) = 1, \quad D(1) = 0.\]
    Therefore, with respect to the bases, we get the desired matrix representation of $D$.
\end{sol}

\begin{exercise}[3c5]
    Suppose $V$ and $W$ are finite-dimensional and $T \in \lin{V, W}$. Prove that there exist a basis of $V$ and a basis of $W$ such that with respect to these bases, all entries of $\mm(T)$ are $0$ except that the entries in row $k$, column $k$, equal $1$ if $1 \leq k \leq \dim \range T$.
\end{exercise}

\begin{sol}
    Let $n = \dim V$ and $m = \dim W$. By the rank-nullity theorem, we have
    \[\dim V = \dim \nul T + \dim \range T.\]
    Let $r = \dim \range T$. Let $u_1, \ldots, u_{n - r}$ be a basis of $\nul T$. Extend this to a basis $u_1, \ldots, u_{n - r}, v_1, \ldots, v_r$ of $V$. We claim that $Tv_1, \ldots, Tv_r$ are linearly independent. Suppose for scalars $\alpha_1, \ldots, \alpha_r \in \f$, we have
    \[\alpha_1 Tv_1 + \cdots + \alpha_r Tv_r = 0.\]
    Then
    \[T(\alpha_1 v_1 + \cdots + \alpha_r v_r) = 0,\]
    which implies that $\alpha_1 v_1 + \cdots + \alpha_r v_r \in \nul T$. Since $u_1, \ldots, u_{n - r}, v_1, \ldots, v_r$ is a basis of $V$, it follows that $\alpha_1 = \cdots = \alpha_r = 0$. Therefore, $Tv_1, \ldots, Tv_r$ are linearly independent and form a basis of $\range T$. Extend this to a basis $Tv_1, \ldots, Tv_r, w_1, \ldots, w_{m - r}$ of $W$.

    If needed, reorder the basis of $V$ so that $v_1, \ldots, v_r$ are the first $r$ basis vectors, and reorder the basis of $W$ so that $Tv_1, \ldots, Tv_r$ are the first $r$ basis vectors. Then with respect to these bases, we have
    \[Tv_k = Tv_k \quad \text{for } k = 1, \ldots, r,\]
    and
    \[Tu_j = 0 \quad \text{for } j = 1, \ldots, n - r.\]
    Therefore, the matrix of $T$ with respect to these bases has $1$'s in the diagonal entries corresponding to $Tv_k$ and $0$'s elsewhere, as desired.
\end{sol}

\begin{exercise}[3c6]
    Suppose $v_1, \ldots, v_m$ is a basis of $V$ and $W$ is finite-dimensional. Suppose $T \in \lin{V, W}$. Prove that there exists a basis $w_1, \ldots, w_n$ of $W$ such that all entries in the first column of $\mm(T)$ [with respect to the bases $v_1, \ldots, v_m$ and $w_1, \ldots, w_n$] are $0$ except for possibly a $1$ in the first row, first column.
\end{exercise}

\begin{sol}
    Let $n = \dim W$. If $Tv_1 = 0$, then we can choose any basis of $W$ and the first column of $\mm(T)$ will be all zeros. If $Tv_1 \neq 0$, then we can extend $Tv_1$ to a basis of $W$. Let $w_1 = Tv_1$ and extend it to a basis $w_1, w_2, \ldots, w_n$ of $W$. Then with respect to the bases $v_1, \ldots, v_m$ of $V$ and $w_1, \ldots, w_n$ of $W$, we have
    \[Tv_1 = w_1\]
    and for each $j = 2, \ldots, m$, we can express
    \[Tv_j = \alpha_{j1} w_1 + \alpha_{j2} w_2 + \cdots + \alpha_{jn} w_n\]
    for some scalars $\alpha_{j1}, \alpha_{j2}, \ldots, \alpha_{jn} \in \f$. Therefore, the first column of $\mm(T)$ has a $1$ in the first row and zeros in all other rows, as desired.
\end{sol}

\begin{exercise}[3c7]
    Suppose $w_1, \ldots, w_n$ is a basis of $W$ and $V$ is finite-dimensional. Suppose $T \in \lin{V, W}$. Prove that there exists a basis $v_1, \ldots, v_m$ of $V$ such that all entries in the first row of $\mm(T)$ [with respect to the bases $v_1, \ldots, v_m$ and $w_1, \ldots, w_n$] are $0$ except for possibly a $1$ in the first row, first column.
\end{exercise}

\begin{sol}
    Let $m = \dim V$ and $u_1, \ldots, u_m$ be a basis for $V$. If $T$ is the zero map or has trivial $w_1$ component, meaning that the coefficient of $w_1$ in $Tv$ is zero for all $v \in V$, then choose any basis of $V$ and the first row of $\mm(T)$ will be all zeros. If at least one $Tu_j$ has non-trivial $w_1$ component, reorder the basis of $V$ so that $u_1$ is such that $Tu_1 \neq 0$. Then we can express
    \[Tu_1 = \alpha w_1 + \beta_2 w_2 + \cdots + \beta_n w_n\]
    for some scalars $\alpha, \beta_2, \ldots, \beta_n \in \f$. Since $Tu_1 \neq 0$, we have $\alpha \neq 0$. We can scale $u_1$ by $\alpha\inv$ to obtain a new vector $v_1 = \alpha\inv u_1$ such that
    \[Tv_1 = w_1 + \frac{\beta_2}{\alpha} w_2 + \cdots + \frac{\beta_n}{\alpha} w_n.\]
    With the remaining $u_2, \ldots, u_m$ basis vectors of $V$, note that they can be expressed as
    \[Tu_j = \alpha_{j1} w_1 + \alpha_{j2} w_2 + \cdots + \alpha_{jn} w_n\]
    for some scalars $\alpha_{j1}, \alpha_{j2}, \ldots, \alpha_{jn} \in \f$. Consider the set of vectors of $V$ given by
    \[v_1 = \alpha\inv u_1, \quad v_j = u_j - \alpha_{j1} v_1 \text{ for } j = 2, \ldots, m.\]
    Note that $Tv_j = Tu_j - \alpha_{j1} Tv_1$ is just a combination of $w_2, \ldots, w_n$. We now claim that $v_1, v_2, \ldots, v_m$ are linearly independent. Suppose for scalars $\gamma_1, \ldots, \gamma_m \in \f$, we have
    \[\gamma_1 v_1 + \gamma_2 v_2 + \cdots + \gamma_m v_m = 0.\]
    Then we can express this as
    \[\gamma_1 v_1 + \gamma_2 (u_2 - \alpha_{21} v_1) + \cdots + \gamma_m (u_m - \alpha_{m1} v_1) = 0.\]
    Simplifying, we get
    \[(\gamma_1 - \gamma_2 \alpha_{21} - \cdots - \gamma_m \alpha_{m1}) v_1 + \gamma_2 u_2 + \cdots + \gamma_m u_m = 0.\]
    Since $v_1, u_2, \ldots, u_m$ are linearly independent, we must have $\gamma_1 - \gamma_2 \alpha_{21} - \cdots - \gamma_m \alpha_{m1} = 0$ and $\gamma_2 = \cdots = \gamma_m = 0$. This implies that $\gamma_1 = 0$, and hence $v_1, v_2, \ldots, v_m$ are linearly independent. Then $v_1, v_2, \ldots, v_m$ form a basis of $V$. It follows that $Tv_1$ has a 1 in the first row, while $Tv_j$ for $j = 2, \ldots, m$ have zeros in the first row.
\end{sol}

\begin{exercise}[3c8]
    Suppose $A$ is an $m$-by-$n$ matrix and $B$ is an $n$-by-$p$ matrix. Prove that $(AB)_{j, \cdot} = A_{j, \cdot} B$ for each $1 \leq j \leq m$. In other words, show that row $j$ of $AB$ equals (row $j$ of $A$) times $B$.
\end{exercise}

\begin{sol}
    Let $A$ be an $m$-by-$n$ matrix and $B$ be an $n$-by-$p$ matrix. The entry in row $j$, column $k$, of the product $AB$ is given by
    \[(AB)_{j, k} = \sum_{r = 1}^{n} A_{j, r} B_{r, k}.\]
    The row vector $A_{j, \cdot}$ is the $j$-th row of $A$, which can be expressed as
    \[A_{j, \cdot} = (A_{j, 1}, A_{j, 2}, \ldots, A_{j, n}).\]
    When we multiply this row vector by the matrix $B$, we get a new row vector whose entries are given by
    \[(A_{j, \cdot} B)_k = \sum_{r = 1}^{n} A_{j, r} B_{r, k}.\]
    Comparing this with the expression for $(AB)_{j, k}$, we see that
    \[(AB)_{j, k} = (A_{j, \cdot} B)_k\]
    for each column index $k$. Therefore, the entire row vector $(AB)_{j, \cdot}$ is equal to the row vector $A_{j, \cdot} B$. This proves that row $j$ of $AB$ equals (row $j$ of $A$) times $B$.
\end{sol}

\begin{exercise}[3c9]
    Suppose $a = (a_1 \quad \cdots \quad a_n)$ is a $1$-by-$n$ matrix and $B$ is an $n$-by-$p$ matrix. Prove that
    \[aB = a_1 B_{1, \cdot} + \cdots + a_n B_{n, \cdot}.\]
    In other words, show that $aB$ is a linear combination of the rows of $B$, with the scalars that multiply the rows coming from $a$.
\end{exercise}

\begin{sol}
    Let $a = (a_1 \quad \cdots \quad a_n)$ be a $1$-by-$n$ matrix and $B$ be an $n$-by-$p$ matrix. The product $aB$ is a $1$-by-$p$ matrix, and its entries are given by
    \[(aB)_{1, k} = \sum_{j = 1}^{n} a_j B_{j, k}\]
    for each column index $k$. 

    Now, consider the linear combination of the rows of $B$ given by
    \[a_1 B_{1, \cdot} + a_2 B_{2, \cdot} + \cdots + a_n B_{n, \cdot}.\]
    The entry in row $1$, column $k$, of this linear combination is
    \[(a_1 B_{1, \cdot} + a_2 B_{2, \cdot} + \cdots + a_n B_{n, \cdot})_k = a_1 B_{1, k} + a_2 B_{2, k} + \cdots + a_n B_{n, k} = \sum_{j = 1}^{n} a_j B_{j, k}.\]

    Comparing this with the expression for $(aB)_{1, k}$, we see that
    \[(aB)_{1, k} = (a_1 B_{1, \cdot} + a_2 B_{2, \cdot} + \cdots + a_n B_{n, \cdot})_k\]
    for each column index $k$. Therefore, we conclude that
    \[aB = a_1 B_{1, \cdot} + a_2 B_{2, \cdot} + \cdots + a_n B_{n, \cdot},\]
    which shows that $aB$ is indeed a linear combination of the rows of $B$, with the scalars coming from the entries of $a$.
\end{sol}

\begin{exercise}[3c10]
    Give an example of $2$-by-$2$ matrices $A$ and $B$ such that $AB \neq BA$.
\end{exercise}

\begin{sol}
    Let 
    \[A =
    \begin{pmatrix}
        1 & 2 \\
        3 & 4
    \end{pmatrix} \quad \text{and} \quad B = 
    \begin{pmatrix}
        0 & 1 \\
        1 & 0
    \end{pmatrix}.\]
    Then we compute the products $AB$ and $BA$:
    \[AB = 
    \begin{pmatrix}
        1 & 2 \\
        3 & 4
    \end{pmatrix} 
    \begin{pmatrix}
        0 & 1 \\
        1 & 0
    \end{pmatrix} 
    = 
    \begin{pmatrix}
        1 \cdot 0 + 2 \cdot 1 & 1 \cdot 1 + 2 \cdot 0 \\
        3 \cdot 0 + 4 \cdot 1 & 3 \cdot 1 + 4 \cdot 0
    \end{pmatrix} 
    = 
    \begin{pmatrix}
        2 & 1 \\
        4 & 3
    \end{pmatrix},\]
    and
    \[BA = 
    \begin{pmatrix}
        0 & 1 \\
        1 & 0
    \end{pmatrix} 
    \begin{pmatrix}
        1 & 2 \\
        3 & 4
    \end{pmatrix} 
    = 
    \begin{pmatrix}
        0 \cdot 1 + 1 \cdot 3 & 0 \cdot 2 + 1 \cdot 4 \\
        1 \cdot 1 + 0 \cdot 3 & 1 \cdot 2 + 0 \cdot 4
    \end{pmatrix} 
    = 
    \begin{pmatrix}
        3 & 4 \\
        1 & 2
    \end{pmatrix}.\]
    By comparing the two products, we see that $AB \neq BA$.
\end{sol}

\begin{exercise}[3c11]
    Prove that the distributive property holds for matrix addition and matrix multiplication. In other words, suppose $A, B, C, D, E$, and $F$ are matrices whose sizes are such that $A(B + C)$ and $(D + E)F$ make sense. Explain why $AB + AC$ and $DF + EF$ both make sense and prove that
    \[A(B + C) = AB + AC \quad \text{and} \quad (D + E)F = DF + EF.\]
\end{exercise}

\begin{sol}
    Let $A$ be an $m$-by-$n$ matrix, and let $B$ and $C$ be $n$-by-$p$ matrices. Then the product $A(B + C)$ is defined. By \autoref{ex:3c8} and \autoref{ex:3c9}, we have
    \[A(B + C)_{j, \cdot} = A_{j, \cdot} (B + C) = A_{j, \cdot} B + A_{j, \cdot} C = (AB)_{j, \cdot} + (AC)_{j, \cdot} = (AB + AC)_{j, \cdot}\]
    for all $j = 1, \ldots, m$. Therefore, $A(B + C) = AB + AC$.

    Similarly, let $D$ and $E$ be $m$-by-$n$ matrices, and let $F$ be an $n$-by-$p$ matrix. Then the product $(D + E)F$ is defined. By \autoref{ex:3c8} and \autoref{ex:3c9}, we have
    \[(D + E)F_{j, \cdot} = (D + E)_{j, \cdot} F = D_{j, \cdot} F + E_{j, \cdot} F = (DF)_{j, \cdot} + (EF)_{j, \cdot} = (DF + EF)_{j, \cdot}\]
    for all $j = 1, \ldots, m$. Therefore, $(D + E)F = DF + EF$.
\end{sol}

\begin{exercise}[3c12]
    Prove that matrix multiplication is associative. In other words, suppose $A$, $B$, and $C$ are matrices whose sizes are such that $(AB)C$ makes sense. Explain why $A(BC)$ makes sense and prove that
    \[(AB)C = A(BC).\]
    \remark{Unfortunately, we cannot use the same trick as in \autoref{ex:3c11} to prove this result.}
\end{exercise}

\begin{sol}
    Let $A$ be an $m$-by-$n$ matrix, $B$ be an $n$-by-$p$ matrix, and $C$ be a $p$-by-$q$ matrix. Then the product $(AB)C$ is defined. Going by the definition of matrix multiplication, we have
    \[((AB)C)_{j, k} = \sum_{r = 1}^{p} (AB)_{j, r} C_{r, k} = \sum_{r = 1}^{p} \left( \sum_{s = 1}^{n} A_{j, s} B_{s, r} \right) C_{r, k} = \sum_{r = 1}^{p} \sum_{s = 1}^{n} A_{j, s} B_{s, r} C_{r, k}.\]
    On the other hand, the product $A(BC)$ is also defined. Again, by the definition of matrix multiplication, we have
    \[(A(BC))_{j, k} = \sum_{s = 1}^{n} A_{j, s} (BC)_{s, k} = \sum_{s = 1}^{n} A_{j, s} \left( \sum_{r = 1}^{p} B_{s, r} C_{r, k} \right) = \sum_{s = 1}^{n} \sum_{r = 1}^{p} A_{j, s} B_{s, r} C_{r, k}.\]
    Since the order of summation does not matter, we can interchange the order of the sums in the last expression to get
    \[(A(BC))_{j, k} = \sum_{r = 1}^{p} \sum_{s = 1}^{n} A_{j, s} B_{s, r} C_{r, k}.\]
    Comparing the expressions for $((AB)C)_{j, k}$ and $(A(BC))_{j, k}$, we see that they are equal for all $j$ and $k$.
\end{sol}

\begin{exercise}[3c13]
    Suppose $A$ is an $n$-by-$n$ matrix and $1 \leq j, k \leq n$. Show that the entry in row $j$, column $k$, of $A^3$ (which is defined to mean $AAA$) is
    \[\sum_{p = 1}^{n} \sum_{r = 1}^{n} A_{j, p} A_{p, r} A_{r, k}.\]
\end{exercise}

\begin{sol}
    Note that $A^3 = (AA)A = A^2 A$. By the definition of matrix multiplication, we have
    \[(A^3)_{j, k} = (A^2 A)_{j, k} = \sum_{p = 1}^{n} (A^2)_{j, p} A_{p, k}.\]
    Now, we can express $(A^2)_{j, p}$ as
    \[(A^2)_{j, p} = (AA)_{j, p} = \sum_{r = 1}^{n} A_{j, r} A_{r, p}.\]
    Substituting this back into the expression for $(A^3)_{j, k}$, we get
    \[(A^3)_{j, k} = \sum_{p = 1}^{n} \left( \sum_{r = 1}^{n} A_{j, r} A_{r, p} \right) A_{p, k} = \sum_{p = 1}^{n} \sum_{r = 1}^{n} A_{j, r} A_{r, p} A_{p, k}.\]
    By renaming the indices $r$ and $p$, we can rewrite this as
    \[(A^3)_{j, k} = \sum_{p = 1}^{n} \sum_{r = 1}^{n} A_{j, p} A_{p, r} A_{r, k},\]
    which is the desired expression.
\end{sol}

\begin{exercise}[3c14]
    Suppose $m$ and $n$ are positive integers. Prove that the function $A \mapsto A^t$ is a linear map from $\f^{m, n}$ to $\f^{n, m}$.
\end{exercise}

\begin{sol}
    Let $A, B \in \f^{m, n}$ and let $\lambda \in \f$. We need to show that $(A + B)^t = A^t + B^t$ and $(\lambda A)^t = \lambda A^t$.

    First, consider the sum $A + B$. The entry in row $j$, column $k$, of $A + B$ is given by
    \[(A + B)_{j, k} = A_{j, k} + B_{j, k}.\]
    Taking the transpose, we have
    \[((A + B)^t)_{k, j} = (A + B)_{j, k} = A_{j, k} + B_{j, k} = (A^t)_{k, j} + (B^t)_{k, j} = (A^t + B^t)_{k, j}.\]
    Therefore, $(A + B)^t = A^t + B^t$.

    Next, consider the scalar multiplication $\lambda A$. The entry in row $j$, column $k$, of $\lambda A$ is given by
    \[(\lambda A)_{j, k} = \lambda A_{j, k}.\]
    Taking the transpose, we have
    \[((\lambda A)^t)_{k, j} = (\lambda A)_{j, k} = \lambda A_{j, k} = \lambda (A^t)_{k, j} = (\lambda A^t)_{k, j}.\]
    Therefore, $(\lambda A)^t = \lambda A^t$.

    Since both properties hold, we conclude that the function $A \mapsto A^t$ is a linear map from $\f^{m, n}$ to $\f^{n, m}$.
\end{sol}

\begin{exercise}[3c15]
    Prove that if $A$ is an $m$-by-$n$ matrix and $C$ is an $n$-by-$p$ matrix, then
    \[(AC)^t = C^t A^t.\]
\end{exercise}

\begin{sol}
    Let $A$ be an $m$-by-$n$ matrix and $C$ be an $n$-by-$p$ matrix. The entry in row $j$, column $k$, of the product $AC$ is given by
    \[(AC)_{j, k} = \sum_{r = 1}^{n} A_{j, r} C_{r, k}.\]
    Taking the transpose, we have
    \[((AC)^t)_{k, j} = (AC)_{j, k} = \sum_{r = 1}^{n} A_{j, r} C_{r, k}.\]

    Now, consider the product $C^t A^t$. The entry in row $k$, column $j$, of this product is given by
    \[(C^t A^t)_{k, j} = \sum_{r = 1}^{n} (C^t)_{k, r} (A^t)_{r, j} = \sum_{r = 1}^{n} C_{r, k} A_{j, r}.\]

    Since multiplication in the field $\f$ is commutative, we can rewrite this as
    \[(C^t A^t)_{k, j} = \sum_{r = 1}^{n} A_{j, r} C_{r, k},\]
    which is exactly the same as the expression for $((AC)^t)_{k, j}$.

    Therefore, we conclude that
    \[(AC)^t = C^t A^t.\]
\end{sol}

\begin{exercise}[3c16]
    Suppose $A$ is an $m$-by-$n$ matrix with $A \neq 0$. Prove that the rank of $A$ is $1$ if and only if there exist $(c_1, \ldots, c_m) \in \f^m$ and $(d_1, \ldots, d_n) \in \f^n$ such that $A_{j, k} = c_j d_k$ for every $j = 1, \ldots, m$ and every $k = 1, \ldots, n$.
\end{exercise}

\begin{sol}
    Suppose the rank of $A$ is 1. Then the column rank of $A$ is 1, which means that all columns of $A$ are scalar multiples of a single non-zero column vector. Let $c = (c_1, \ldots, c_m)^t$ be this non-zero column vector. Then for each column $k$ of $A$, there exists a scalar $d_k \in \f$ such that the $k$-th column of $A$ is given by $d_k c$. Therefore, we can express the entries of $A$ as
    \[A_{j, k} = c_j d_k\]
    for all $j = 1, \ldots, m$ and $k = 1, \ldots, n$.

    Conversely, suppose there exist $(c_1, \ldots, c_m) \in \f^m$ and $(d_1, \ldots, d_n) \in \f^n$ such that $A_{j, k} = c_j d_k$ for all $j$ and $k$. Since $A \neq 0$, then $(c_1, \ldots, c_m) \neq 0$ and $(d_1, \ldots, d_n) \neq 0$. This means that all columns of $A$ are scalar multiples of the column vector $c$, and hence the column rank of $A$ is 1. Therefore, the rank of $A$ is 1.
\end{sol}

\begin{exercise}[3c17]
    Suppose $T \in \lin{V}$, and $u_1, \ldots, u_n$ and $v_1, \ldots, v_n$ are bases of $V$. Prove that the following are equivalent.
    \begin{subexercises}
        \item $T$ is injective.
        \item The columns of $\mm(T)$ are linearly independent in $\f^{n, 1}$.
        \item The columns of $\mm(T)$ span $\f^{n, 1}$.
        \item The rows of $\mm(T)$ span $\f^{1, n}$.
        \item The rows of $\mm(T)$ are linearly independent in $\f^{1, n}$.
    \end{subexercises}
\end{exercise}

\begin{solitem}
    \item (a) $\iff$ (b): (a) $\iff \dim \nul T = 0 \iff \dim \range T = n \iff \colrank \mm(T) = n$. For $n$ column vectors, $\rank n$ implies linear independence $\iff$ (b).
    \item (b) $\iff$ (c): $n$ columns of $\mm(T)$ are linearly independent in $\f^{n, 1} \iff n$ columns of $\mm(T)$ span $\f^{n, 1}$.
    \item (c) $\iff$ (d): columns span $\f^{n, 1} \iff \rank \mm(T) = n$ by (b) $\iff$ (c). By Theorem 3.57, $\rank \mm(T) = \rowrank \mm(T) \iff$ rows span $\f^{1, n}$.
    \item (d) $\iff$ (e): $n$ rows of $\mm(T)$ span $\f^{1, n} \iff n$ rows of $\mm(T)$ are linearly independent in $\f^{1, n}$.
\end{solitem}